\documentclass[11pt]{article}
\usepackage[a4paper,margin=22mm]{geometry}
\usepackage{fontspec}
\setmainfont{DejaVu Serif}
\setsansfont{DejaVu Sans}
\usepackage{microtype}
\usepackage{xcolor}
\usepackage{hyperref}
\usepackage{enumitem}
\usepackage{parskip}
\definecolor{Accent}{HTML}{971D32}
\definecolor{Muted}{HTML}{666666}
\hypersetup{colorlinks=true,urlcolor=Accent,linkcolor=Accent}
\setlist{leftmargin=1.6em,itemsep=0.3em,topsep=0.35em}
\setcounter{tocdepth}{2}
\renewcommand{\contentsname}{Contents}
\newcommand{\sectionrule}{\par\vspace{-0.5em}\noindent\textcolor{Accent}{\rule{\linewidth}{0.6pt}}\par}
\begin{document}

\begin{center}
{\sffamily\bfseries\color{Accent} TOEFL WORD OF THE DAY\par}
\vspace{8mm}
{\fontsize{42}{48}\selectfont\bfseries deflect\par}
\vspace{2mm}
{\Large\sffamily\color{Accent} /dɪˈflɛkt/\par}
\vspace{2mm}
{\small\sffamily Local audio phrase: deflect\par}
{\small\sffamily Audio file: audios/deflect.mp3\par}
\vspace{2mm}
{\large\sffamily\color{Muted} transitive and intransitive verb\par}
\vspace{5mm}
{\sffamily cause something to change direction \textbullet\ change direction \textbullet\ redirect attention or criticism\par}
\vspace{6mm}
{\large To deflect something is to turn it away from its original direction, either physically or figuratively.\par}
\vspace{6mm}
\begin{minipage}{0.84\textwidth}
\itshape “Earth's magnetic field helps deflect charged particles from the Sun.”
\end{minipage}\par
\vspace{10mm}
\textbf{Date:} August 12th 2026\quad\textbullet\quad \textbf{Level:} B1--mid-B2 TOEFL
\vfill
Created and compiled by \textbf{Amir Kooshky}\par
\vspace{2mm}
\href{https://t.me/KooshkyTOEFL}{Telegram Channel}\quad\textbullet\quad
\href{https://www.instagram.com/kooshkytoefl}{Instagram}\quad\textbullet\quad
\href{https://t.me/KooshkyTOEFL\_pv}{Personal Telegram}
\end{center}
\newpage
\tableofcontents
\newpage

\section{Meanings and usage}
\sectionrule
\subsection{Meaning 1: Cause to change direction}
\textbf{Part of speech:} transitive verb

\textbf{IPA:} /dɪˈflɛkt/

\textbf{Pronunciation phrase:} deflect

\textbf{Local audio file:} audios/deflect.mp3

\textbf{Register:} neutral to formal; common in science, engineering, physics, and technical descriptions

\textbf{Definition:} to cause something that is moving to change direction

\textbf{In simpler words:} make something turn away

\textbf{Typical pattern:} deflect + object + away from or toward + direction

\subsubsection{Natural examples}
\begin{enumerate}
  \item The protective shield is designed to deflect small pieces of flying debris.
  \item Earth's magnetic field helps deflect charged particles from the Sun.
  \item The angled surface deflects wind away from the entrance.
  \item The barrier can deflect waves before they reach the harbor.
  \item Engineers added a panel to deflect heat away from sensitive equipment.
  \item The curved structure helps deflect rainwater toward the drainage system.
  \item The material was strong enough to deflect some of the impact.
\end{enumerate}

\subsubsection{Nuance and implication}
\textbf{Central nuance:} The object continues moving, but its path is changed by a force, surface, barrier, or other influence.

\textbf{Usual implication:} The change may protect something, redirect energy, or alter where a moving object eventually goes.

\textbf{Compared with stop:} To stop something is to end its movement. To deflect it is to change its direction while it continues moving.

\subsubsection{Grammar notes}
\textbf{How to use it:} This verb normally needs the moving object or force after it.

\textbf{What can follow it:} It is commonly followed by nouns such as light, radiation, wind, water, debris, particles, heat, or an object in motion.

\textbf{Common preposition:} Use away from when showing what the object is being redirected away from, as in deflect heat away from the engine.

\textbf{Position in a sentence:} The thing causing the change normally comes before the verb, and the moving object comes directly after it.

\subsubsection{Useful patterns}
\textbf{Pattern:} deflect + object

\textbf{Use:} Use this when something causes a moving object or force to change direction.

\textbf{Example:} The shield deflected the incoming particles.

\textbf{Pattern:} deflect + object + away from + noun

\textbf{Use:} Use this when showing what is being protected from the moving object or force.

\textbf{Example:} The device deflects heat away from the battery.

\textbf{Pattern:} be designed to deflect + noun

\textbf{Use:} Use this for equipment or structures created to redirect something.

\textbf{Example:} The wall was designed to deflect strong winds.

\subsubsection{Common wrong usage}
\paragraph{Correction 1}
\textbf{Wrong:} The shield deflected away the particles.

\textbf{Correct:} The shield deflected the particles away.

\textbf{Why:} Place the object directly after deflect.

\paragraph{Correction 2}
\textbf{Wrong:} The barrier deflected the water to hit the building.

\textbf{Correct:} The barrier deflected the water away from the building.

\textbf{Why:} Use away from when the purpose is to redirect something so that it does not hit a particular place.

\subsection{Meaning 2: Change direction}
\textbf{Part of speech:} intransitive verb

\textbf{IPA:} /dɪˈflɛkt/

\textbf{Pronunciation phrase:} deflect

\textbf{Local audio file:} audios/deflect.mp3

\textbf{Register:} neutral to formal; common in physical descriptions and technical contexts

\textbf{Definition:} to change direction after hitting something or being affected by a force

\textbf{In simpler words:} turn in a different direction

\textbf{Typical pattern:} deflect off + surface or deflect to + direction

\subsubsection{Natural examples}
\begin{enumerate}
  \item The ball deflected off a defender and entered the goal.
  \item Some of the incoming radiation is deflected by particles in the atmosphere.
  \item The stream of water deflected to the side when it struck the metal plate.
  \item The projectile deflected upward after striking the hard surface.
  \item Light can deflect as it passes through certain materials.
  \item The ball deflected off the post and rolled across the field.
  \item The beam deflected slightly when the equipment was moved.
\end{enumerate}

\subsubsection{Nuance and implication}
\textbf{Central nuance:} Here the subject itself changes direction rather than causing another thing to change direction.

\textbf{Usual implication:} The change often happens because the moving object strikes a surface or encounters another force.

\subsubsection{Grammar notes}
\textbf{How to use it:} In this meaning, the moving object itself is the subject, so the verb does not need an object.

\textbf{What can follow it:} It is often followed by off and the surface that was hit, or by a word showing the new direction.

\textbf{Common preposition:} Use deflect off something when the object changes direction after contact, as in The ball deflected off the wall.

\subsubsection{Useful patterns}
\textbf{Pattern:} deflect off + surface or object

\textbf{Use:} Use this when something changes direction after hitting another object.

\textbf{Example:} The ball deflected off the goalkeeper.

\textbf{Pattern:} deflect to + direction

\textbf{Use:} Use this to state the direction in which something turns.

\textbf{Example:} The stream deflected to the left.

\textbf{Pattern:} be deflected by + force or object

\textbf{Use:} Use the passive form when emphasizing what caused the change in direction.

\textbf{Example:} The particles were deflected by the magnetic field.

\subsubsection{Common wrong usage}
\paragraph{Correction 1}
\textbf{Wrong:} The ball deflected from the defender and changed direction.

\textbf{Correct:} The ball deflected off the defender and changed direction.

\textbf{Why:} Use deflect off when something changes direction after hitting another object.

\paragraph{Correction 2}
\textbf{Wrong:} The beam was deflected from the magnetic field.

\textbf{Correct:} The beam was deflected by the magnetic field.

\textbf{Why:} Use by to identify the force or object that causes the deflection.

\subsection{Meaning 3: Redirect attention or criticism}
\textbf{Part of speech:} transitive verb

\textbf{IPA:} /dɪˈflɛkt/

\textbf{Pronunciation phrase:} deflect

\textbf{Local audio file:} audios/deflect.mp3

\textbf{Register:} neutral to formal; common in politics, journalism, business, and discussions of communication

\textbf{Definition:} to turn attention, criticism, blame, or a difficult question away from a person or subject

\textbf{In simpler words:} turn attention or blame away

\textbf{Typical pattern:} deflect + criticism, blame, attention, question, or responsibility

\subsubsection{Natural examples}
\begin{enumerate}
  \item The spokesperson tried to deflect criticism by emphasizing the company's recent achievements.
  \item The minister deflected questions about the failed project and discussed future plans instead.
  \item Some leaders attempt to deflect blame onto other departments when problems arise.
  \item The company tried to deflect attention away from the environmental impact of the project.
  \item Rather than answering directly, he deflected the question with a joke.
  \item Officials accused the organization of trying to deflect responsibility for the accident.
  \item The advertisement was intended to deflect public attention from the recent controversy.
\end{enumerate}

\subsubsection{Nuance and implication}
\textbf{Central nuance:} This figurative use treats attention, blame, or criticism as if it were something moving toward a person and being turned elsewhere.

\textbf{Usual implication:} It often suggests avoiding responsibility, an uncomfortable topic, or a direct answer.

\textbf{Compared with address:} To address criticism or a question is to deal with it directly. To deflect it is to redirect attention away from it.

\subsubsection{Grammar notes}
\textbf{How to use it:} This verb normally needs the criticism, blame, attention, question, or responsibility after it.

\textbf{What can follow it:} Common objects include criticism, blame, attention, questions, responsibility, suspicion, and pressure.

\textbf{Common preposition:} Use away from to show what attention is being redirected from. With blame, onto is common when naming the new target, as in deflect blame onto someone else.

\textbf{Position in a sentence:} It is often followed by by and an -ing form to explain how the deflection happens.

\subsubsection{Useful patterns}
\textbf{Pattern:} deflect criticism

\textbf{Use:} Use this when someone tries to reduce or redirect criticism.

\textbf{Example:} The company released new figures in an attempt to deflect criticism.

\textbf{Pattern:} deflect attention away from + noun

\textbf{Use:} Use this when someone tries to make people focus on something else.

\textbf{Example:} The announcement deflected attention away from the budget problems.

\textbf{Pattern:} deflect blame onto + person or group

\textbf{Use:} Use this when someone tries to make another person appear responsible.

\textbf{Example:} The manager tried to deflect blame onto junior employees.

\textbf{Pattern:} deflect a question

\textbf{Use:} Use this when someone avoids answering a question directly.

\textbf{Example:} The politician deflected the question and changed the subject.

\textbf{Pattern:} deflect something by + -ing form

\textbf{Use:} Use this to explain the method used to redirect criticism or attention.

\textbf{Example:} The company tried to deflect criticism by announcing a new environmental program.

\subsubsection{Common wrong usage}
\paragraph{Correction 1}
\textbf{Wrong:} The minister deflected from the question.

\textbf{Correct:} The minister deflected the question.

\textbf{Why:} When deflect means avoid or redirect a question, place the question directly after the verb.

\paragraph{Correction 2}
\textbf{Wrong:} The company deflected blame to its suppliers.

\textbf{Correct:} The company deflected blame onto its suppliers.

\textbf{Why:} Use onto when blame is redirected toward another person or group.

\paragraph{Correction 3}
\textbf{Wrong:} She deflected attention from herself to not answer.

\textbf{Correct:} She deflected attention away from herself to avoid answering.

\textbf{Why:} Use away from for the original focus and to avoid when explaining the purpose.

\section{Word family}
\sectionrule
\subsection{deflection}
\textbf{Part of speech:} countable or uncountable noun

\textbf{IPA:} /dɪˈflɛkʃən/

\textbf{Definition:} a change in direction, or an attempt to redirect attention, criticism, or responsibility

\textbf{Relationship to the headword:} This noun names the physical change in direction or the figurative act of redirecting attention or criticism.

\textbf{Grammar pattern:} the deflection of + object or a deflection from + issue

\textbf{Usage note:} It is common in physics and engineering, and it is also used figuratively in discussions of communication.

\subsubsection{Examples}
\begin{enumerate}
  \item The magnetic field caused a slight deflection of the particles.
  \item Critics described his response as a deflection rather than a direct answer.
  \item Engineers measured the deflection of the beam under pressure.
\end{enumerate}

\subsection{deflector}
\textbf{Part of speech:} countable noun

\textbf{IPA:} /dɪˈflɛktər/

\textbf{Definition:} a device or surface designed to change the direction of something such as air, heat, water, or particles

\textbf{Relationship to the headword:} A deflector is something designed to deflect a moving substance or force.

\textbf{Grammar pattern:} air, wind, heat, or water + deflector

\textbf{Usage note:} It is mainly used in technical and engineering contexts.

\subsubsection{Examples}
\begin{enumerate}
  \item The machine uses a metal deflector to direct hot air away from the motor.
  \item A wind deflector was installed above the entrance.
\end{enumerate}

\section{Common collocations}
\sectionrule
\subsection{deflect attention}
\textbf{Applies to:} Redirect attention or criticism

\textbf{Meaning:} cause people to focus on something else

\textbf{Pattern:} deflect attention away from + issue

\textbf{Example:} The announcement was intended to deflect attention from the investigation.

\subsection{deflect criticism}
\textbf{Applies to:} Redirect attention or criticism

\textbf{Meaning:} redirect or reduce criticism directed at someone

\textbf{Example:} Officials attempted to deflect criticism by promising an independent review.

\subsection{deflect blame}
\textbf{Applies to:} Redirect attention or criticism

\textbf{Meaning:} redirect responsibility for a problem toward someone else

\textbf{Pattern:} deflect blame onto + person or group

\textbf{Example:} The executive tried to deflect blame onto the previous management team.

\subsection{deflect a question}
\textbf{Applies to:} Redirect attention or criticism

\textbf{Meaning:} avoid answering a question directly

\textbf{Example:} The spokesperson deflected the question and returned to the prepared statement.

\subsection{deflect responsibility}
\textbf{Applies to:} Redirect attention or criticism

\textbf{Meaning:} avoid accepting responsibility by directing it elsewhere

\textbf{Example:} The report accused senior managers of trying to deflect responsibility for the failure.

\subsection{deflect heat}
\textbf{Applies to:} Cause to change direction

\textbf{Meaning:} redirect heat away from a particular place

\textbf{Pattern:} deflect heat away from + object

\textbf{Example:} The reflective surface helps deflect heat away from the building.

\subsection{deflect radiation}
\textbf{Applies to:} Cause to change direction

\textbf{Meaning:} cause radiation or charged particles to move in another direction

\textbf{Example:} The magnetic field can deflect some forms of charged radiation.

\subsection{deflect off}
\textbf{Applies to:} Change direction

\textbf{Meaning:} change direction after hitting something

\textbf{Pattern:} deflect off + surface or object

\textbf{Example:} The ball deflected off the wall and returned toward the player.

\subsection{be deflected by}
\textbf{Applies to:} Change direction

\textbf{Meaning:} have one's direction changed by a force or object

\textbf{Pattern:} be deflected by + force or object

\textbf{Example:} The particles were deflected by Earth's magnetic field.

\subsection{slight deflection}
\textbf{Applies to:} Change direction

\textbf{Meaning:} a small change in direction

\textbf{Example:} The experiment produced a slight deflection in the path of the beam.

\textbf{Usage note:} Deflection is the noun form.

\section{Synonyms and antonyms}
\sectionrule
\subsection{Close synonyms}
\subsubsection{redirect}
\textbf{Part of speech:} transitive verb

\textbf{Applies to:} all senses

\textbf{Shared meaning:} Both verbs can mean cause something to move or focus in a different direction.

\textbf{Difference:} Redirect is a broad, neutral word for sending something in a different direction. Deflect often suggests that the change happens because of contact, force, resistance, or an attempt to avoid something.

\textbf{Example:} The system redirects excess water into a storage tank.

\subsubsection{divert}
\textbf{Part of speech:} transitive verb

\textbf{Applies to:} all senses

\textbf{Shared meaning:} Both verbs can describe changing direction or shifting attention.

\textbf{Difference:} Divert often means deliberately send something along a different route or shift resources or attention elsewhere. Deflect often suggests turning something aside from its original path or target.

\textbf{Example:} Traffic was diverted around the damaged bridge.

\subsubsection{avoid}
\textbf{Part of speech:} transitive verb

\textbf{Applies to:} Redirect attention or criticism

\textbf{Shared meaning:} Both can describe not dealing directly with an uncomfortable issue.

\textbf{Difference:} Avoid means stay away from or not deal with something. Deflect is more specific and suggests redirecting the question, criticism, or attention elsewhere.

\textbf{Example:} The official avoided answering the most difficult question.

\subsection{Direct or useful antonyms}
\subsubsection{address}
\textbf{Part of speech:} transitive verb

\textbf{Applies to:} Redirect attention or criticism

\textbf{Opposition:} To address a question or criticism is to deal with it directly, while to deflect it is to redirect attention away from it.

\textbf{Example:} The minister directly addressed concerns about the new policy.

\section{Words learners may confuse}
\sectionrule
\subsection{reflect}
\textbf{Part of speech:} verb

\textbf{Why learners confuse them:} Both verbs can describe a change in the path of energy or moving particles.

\textbf{Key difference:} Deflect means change the direction of something. Reflect can mean send light, sound, or heat back from a surface, often toward the side it came from.

\textbf{deflect:} The shield deflects charged particles away from the spacecraft.

\textbf{reflect:} The mirror reflects light back toward the observer.

\subsection{distract}
\textbf{Part of speech:} transitive verb

\textbf{Why learners confuse them:} Both can involve changing what someone is paying attention to.

\textbf{Key difference:} Deflect attention means deliberately redirect attention away from a subject. Distract means make it difficult for someone to concentrate, often without giving the attention a specific new target.

\textbf{deflect:} The spokesperson tried to deflect attention from the scandal.

\textbf{distract:} Noise from the hallway distracted the students during the test.

\section{Etymology}
\sectionrule
\textbf{Origin:} Deflect comes from Latin deflectere, meaning to bend or turn aside.

\section{Practice exercises}
\sectionrule
\subsection{Word family choice}
Choose the best member of the word family for each blank.

\begin{enumerate}
  \item Engineers measured the \_\_\_\_\_ of the metal beam under heavy pressure.
  \item The machine includes a heat \_\_\_\_\_ that protects sensitive components.
  \item The shield helps \_\_\_\_\_ small particles away from the spacecraft.
\end{enumerate}

\subsection{Correct or incorrect?}
Decide whether each sentence is correct. Correct the inaccurate ones.

\begin{enumerate}
  \item The barrier deflects wind away from the building.
  \item The ball deflected from the wall and changed direction.
  \item The politician tried to deflect the question by discussing another issue.
  \item The manager deflected blame to junior employees.
\end{enumerate}

\subsection{Collocation practice}
Complete each sentence with a natural collocation.

\begin{enumerate}
  \item The reflective surface helps deflect heat \_\_\_\_\_ from the equipment.
  \item The spokesperson attempted to deflect \_\_\_\_\_ by blaming outside contractors.
  \item The ball deflected \_\_\_\_\_ the defender and changed direction.
  \item The company tried to deflect blame \_\_\_\_\_ its suppliers.
\end{enumerate}

\subsection{Complete answer key}
\subsubsection{Word family choice}
\begin{enumerate}
  \item \textbf{deflection.} Engineers measured the deflection of the metal beam under heavy pressure. \emph{Use the noun deflection for a change in position or direction.}
  \item \textbf{deflector.} The machine includes a heat deflector that protects sensitive components. \emph{Use deflector for a device designed to redirect heat.}
  \item \textbf{deflect.} The shield helps deflect small particles away from the spacecraft. \emph{Use the base verb deflect after helps.}
\end{enumerate}

\subsubsection{Correct or incorrect?}
\begin{enumerate}
  \item \textbf{Correct} \emph{Deflect something away from something is a natural physical pattern.}
  \item \textbf{Incorrect}. The ball deflected off the wall and changed direction. \emph{Use deflect off when something changes direction after hitting a surface.}
  \item \textbf{Correct} \emph{Deflect a question correctly means avoid answering it directly by redirecting attention.}
  \item \textbf{Incorrect}. The manager deflected blame onto junior employees. \emph{Use onto when blame is redirected toward another person or group.}
\end{enumerate}

\subsubsection{Collocation practice}
\begin{enumerate}
  \item \textbf{away.} The reflective surface helps deflect heat away from the equipment. \emph{Deflect something away from something means redirect it so that it does not reach the target.}
  \item \textbf{criticism.} The spokesperson attempted to deflect criticism by blaming outside contractors. \emph{Deflect criticism means redirect or reduce criticism aimed at someone.}
  \item \textbf{off.} The ball deflected off the defender and changed direction. \emph{Use deflect off when a moving object changes direction after contact.}
  \item \textbf{onto.} The company tried to deflect blame onto its suppliers. \emph{Deflect blame onto someone means redirect responsibility toward that person or group.}
\end{enumerate}

\vfill
\begin{center}
\small\color{Muted} Created and compiled by \textbf{Amir Kooshky}\\
\href{https://t.me/KooshkyTOEFL}{Telegram Channel}\quad\textbullet\quad
\href{https://www.instagram.com/kooshkytoefl}{Instagram}\quad\textbullet\quad
\href{https://t.me/KooshkyTOEFL\_pv}{Personal Telegram}
\end{center}
\end{document}
